\section{Robust branching and discontinuity barriers}

\subsection{Overlapping branches}

Exact comparison of a represented real with zero is generally discontinuous.
A robust program replaces a partition by an overlap.  Suppose a rational
approximation \(a\) satisfies \(\abs{x-a}\le\eta\).  Define:
\[
  L(a): a\le \eta,\qquad
  R(a): a\ge-\eta.
\]
At least one branch holds, and both may hold in the overlap
\([-\eta,\eta]\).  A node contract must prove that either branch gives an
acceptable output when its guard holds.  The constructor may choose the first
passing guard by bounded Boolean search.

\begin{lemma}[Overlapping-branch soundness]
If \(\abs{x-a}\le\eta\), \(L(a)\) implies \(x\le2\eta\), and \(R(a)\) implies
\(x\ge-2\eta\).  Moreover \(L(a)\vee R(a)\) is decidable and always true.
\end{lemma}

\begin{proof}
The first two claims are triangle inequalities.  For coverage, if
\(a\le\eta\) then \(L\) holds.  Otherwise the decidable rational inequality
gives \(a>\eta\ge-\eta\), so \(R\) holds.  This is a case split on an exact
rational Boolean.
\end{proof}

This suffices whenever the two branch outputs agree up to the allowed error
on the overlap.  It is the constructive analogue of a stable conditional.

\subsection{Exact decision from a gap}

\begin{theorem}[Gap-decision theorem]
Let \(x\) have a rational approximation \(a\) with
\(\abs{x-a}<\delta/4\), where \(\delta>0\).  Assume the finite evidence
\[
           x=0\quad\vee\quad \abs{x}\ge\delta.
\]
Then the rational test \(\abs a<\delta/2\) decides which alternative holds.
\end{theorem}

\begin{proof}
If \(\abs a<\delta/2\), the apart alternative would imply
\[
 \delta\le\abs x\le\abs{x-a}+\abs a
 <\delta/4+\delta/2<\delta,
\]
a contradiction; hence \(x=0\).  If \(\abs a\ge\delta/2\), then \(x=0\)
would give \(\abs a=\abs{a-x}<\delta/4\), a contradiction; hence
\(\abs x\ge\delta\).  The proof uses the supplied disjunction, not an
unrestricted law of excluded middle.
\end{proof}

The same proof decides among finitely many regions separated by a rational
margin.  It is the basis for exact pivot selection, inversion, square-free
decomposition with a discriminant bound, exact chart choice, and branch
collapse.

\subsection{Why the gap is not optional}

\begin{obstruction}[Uniform zero decision]
There is no total computable procedure that decides whether an arbitrary fast
Cauchy name of a nonnegative real represents zero.
\end{obstruction}

\begin{proof}
Let \(H(e,t)\) be the primitive-recursive Boolean saying that machine \(e\)
halts within \(t\) steps, and define
\[
  A(e,s)=\sum_{t<s}2^{-t-1}H(e,t).
\]
The omitted tail is at most \(2^{-s}\), so \(A(e,\cdot)\) is one fixed
primitive-recursive family of fast Cauchy names.  If the machine never halts,
the limit is zero.  If it first halts at stage \(h\), then \(H(e,t)=1\) for
all \(t\ge h\), and the limit is \(2^{-h}>0\).  A uniform zero decider would
therefore decide the halting problem.
\end{proof}

This metatheoretic statement does not prove that every first-order sentence
called ``root selection'' is unprovable in iPRA.  It identifies a particular
impossible interface: individually accessible exact choices from arbitrary
Cauchy names.  Approximate sets, finite covers, projective outputs, and
gap-conditioned selectors are different and often valid.

\subsection{Projective and quotient outputs}

Some apparent choices are artifacts of representation.  A rotation has two
unit-quaternion representatives \(q,-q\).  An eigenline has many nonzero
vectors and two unit vectors of opposite sign.  A polynomial's unordered
root multiset has many orderings.  The primary output should therefore be a
quotient-invariant certificate:
\[
  qq^\ast=1,\qquad R=R(q),
\]
or a rank-one projector, or a symmetric list code.  Normalization and sign
selection are secondary operations guarded by nonzero and orientation
evidence.

This move is not merely philosophical.  Quotient-invariant polynomial
identities are exact finite arithmetic and compose well in iPRA, while a
canonical representative may force a discontinuous branch.
